%==============================================================================
\section{Mathematics --- Toán học nền tảng cho học máy}
\label{sec:toan-nen-ml}
%==============================================================================

\begin{ynghia}[title={Nguồn}]
\tsurl{https://www.tutorialspoint.com/machine_learning/machine_learning_mathematics.htm}.
\end{ynghia}

\begin{dongco}[title={Vai trò của toán trong ML}]
Học máy là lĩnh vực liên ngành gồm khoa học máy tính, thống kê và toán học.
Toán đóng vai trò then chốt trong việc phát triển và hiểu các thuật toán học máy.
Bài này trình bày các khái niệm toán cần thiết: đại số tuyến tính, giải tích, xác suất và thống kê.
\end{dongco}

\begin{luuy}[title={Toán chi tiết trong repo}]
Phần toán đầy đủ hơn (tensor, sigmoid, gradient, hồi quy logistic, MLP): \path{toan\_dai\_hoc/nhom\_9\_khac/O3\_hoc\_may/ly\_thuyet.tex}, mục ``Toán nền cho học máy''.
\end{luuy}

\begin{luuy}[title={Đối chiếu với bài Khái niệm nền}]
Định nghĩa học máy, train/test và các loại học đã được trình bày ở phần \textbf{Khái niệm nền}; bài này tập trung vào \textbf{nền toán} phục vụ đọc hiểu thuật toán.
\end{luuy}

\subsection{Đại số tuyến tính (Linear Algebra)}

\begin{dinhnghia}[title={Đại số tuyến tính}]
Đại số tuyến tính là nhánh toán học nghiên cứu phương trình tuyến tính và cách biểu diễn chúng trong không gian vector.
Trong học máy, đại số tuyến tính dùng để biểu diễn và thao tác dữ liệu: vector và ma trận biểu diễn điểm dữ liệu, đặc trưng và trọng số trong mô hình.
\end{dinhnghia}

\begin{ynghia}[title={Vector và ma trận}]
Vector là danh sách có thứ tự các số; ma trận là mảng hình chữ nhật các số.
Ví dụ: một vector có thể biểu diễn một điểm dữ liệu; một ma trận có thể biểu diễn cả tập dữ liệu.
Các phép toán tuyến tính (nhân ma trận, nghịch đảo, \ldots) dùng để biến đổi và phân tích dữ liệu.
\end{ynghia}

\begin{phuongphap}[title={Các khái niệm đại số tuyến tính quan trọng trong ML}]
\begin{itemize}
  \item \textbf{Vector và ma trận}: biểu diễn dataset, feature, giá trị mục tiêu, trọng số, \ldots
  \item \textbf{Phép toán ma trận}: cộng, trừ, nhân, chuyển vị --- xuất hiện trong hầu hết thuật toán ML.
  \item \textbf{Giá trị riêng và vector riêng}: hữu ích cho giảm chiều, ví dụ PCA.
  \item \textbf{Chiếu (projection)}: siêu phẳng và chiếu lên mặt phẳng --- cần để hiểu SVM.
  \item \textbf{Phân tích ma trận}: phân tích ma trận, SVD --- trích thông tin quan trọng từ dữ liệu.
  \item \textbf{Tensor}: trong deep learning biểu diễn dữ liệu đa chiều; tensor có thể là vô hướng, vector hoặc ma trận.
  \item \textbf{Gradient}: tìm giá trị tối ưu của tham số mô hình.
  \item \textbf{Ma trận Jacobian}: phân tích quan hệ giữa biến đầu vào và đầu ra của mô hình.
  \item \textbf{Trực giao (orthogonality)}: khái niệm cốt lõi trong PCA, SVM.
\end{itemize}
\end{phuongphap}

\subsection{Giải tích (Calculus)}

\begin{dinhnghia}[title={Giải tích}]
Giải tích là nhánh toán học nghiên cứu tốc độ thay đổi và tích lũy.
Trong ML, giải tích dùng để tối ưu mô hình bằng cách tìm cực tiểu hoặc cực đại của hàm số --- điển hình là thuật toán \textbf{gradient descent}.
\end{dinhnghia}

\begin{phuongphap}[title={Gradient descent (tóm tắt theo nguồn)}]
Gradient descent là thuật toán tối ưu lặp: cập nhật trọng số mô hình theo gradient của hàm mất mát (loss).
Gradient là vector các đạo hàm riêng của loss theo từng trọng số.
Lặp lại cập nhật theo hướng gradient âm nhằm \textbf{giảm} loss.
\end{phuongphap}

\begin{phuongphap}[title={Các khái niệm giải tích quan trọng trong ML}]
\begin{itemize}
  \item \textbf{Hàm số}: cốt lõi của ML --- mô hình học một hàm ánh xạ đầu vào $\rightarrow$ đầu ra khi huấn luyện; cần nắm hàm liên tục/rời rạc.
  \item \textbf{Đạo hàm, gradient, độ dốc}: hiểu cách thuật toán tối ưu (gradient descent) hoạt động.
  \item \textbf{Đạo hàm riêng}: tìm cực đại/cực tiểu; thường dùng trong tối ưu.
  \item \textbf{Quy tắc chuỗi (chain rule)}: tính đạo hàm loss có nhiều biến --- ứng dụng chính trong mạng nơ-ron.
  \item \textbf{Phương pháp tối ưu}: tìm tham số làm cost function nhỏ nhất; gradient descent là phương pháp phổ biến nhất.
\end{itemize}
\end{phuongphap}

\subsection{Lý thuyết xác suất (Probability Theory)}

\begin{dinhnghia}[title={Xác suất}]
Xác suất nghiên cứu sự không chắc chắn và tính ngẫu nhiên.
Trong ML, xác suất mô hình hoá và phân tích dữ liệu không chắc chắn hoặc biến thiên; phân phối xác suất (Gaussian, Poisson, \ldots) mô tả xác suất điểm dữ liệu hoặc sự kiện.
\end{dinhnghia}

\begin{ynghia}[title={Suy luận Bayes}]
Suy luận Bayes là kỹ thuật mô hình hoá xác suất phổ biến trong ML, dựa trên định lý Bayes:
xác suất giả thuyết cho dữ liệu quan sát tỷ lệ với xác suất dữ liệu cho giả thuyết nhân với xác suất tiên nghiệm của giả thuyết.
Cập nhật tiên nghiệm theo dữ liệu quan sát cho phép dự đoán hoặc phân loại theo xác suất.
\end{ynghia}

\begin{phuongphap}[title={Các khái niệm xác suất quan trọng trong ML}]
\begin{itemize}
  \item \textbf{Xác suất đơn giản}: nền tảng; mọi bài toán phân loại đều dùng xác suất; hàm SoftMax trong ANN dùng xác suất đơn giản.
  \item \textbf{Xác suất có điều kiện}: ví dụ bộ phân loại Naive Bayes.
  \item \textbf{Biến ngẫu nhiên}: gán giá trị khởi tạo tham số mô hình --- coi là bước đầu của quá trình huấn luyện.
  \item \textbf{Phân phối xác suất}: dùng khi xây hàm mất mát cho phân loại.
  \item \textbf{Phân phối liên tục và rời rạc}: mô hình hoá các kiểu dữ liệu khác nhau trong ML.
  \item \textbf{Hàm phân phối}: mô hình hoá phân phối sai số trong hồi quy tuyến tính và mô hình thống kê khác.
  \item \textbf{Ước lượng hợp lý cực đại (MLE)}: nền của một số phương pháp ML và deep learning cho phân loại.
\end{itemize}
\end{phuongphap}

\subsection{Thống kê (Statistics)}

\begin{dinhnghia}[title={Thống kê}]
Thống kê nghiên cứu thu thập, phân tích, diễn giải và trình bày dữ liệu.
Trong ML, thống kê dùng để đánh giá/so sánh mô hình, ước lượng tham số và kiểm định giả thuyết.
\end{dinhnghia}

\begin{vidu}[title={Cross-validation}]
Cross-validation là kỹ thuật thống kê đánh giá hiệu năng mô hình trên dữ liệu mới, chưa thấy.
Tập dữ liệu được chia thành nhiều phần; mô hình được huấn luyện và đánh giá trên từng phần, từ đó ước lượng hiệu năng trên dữ liệu mới và so sánh các mô hình.
\end{vidu}

\begin{phuongphap}[title={Các khái niệm thống kê quan trọng trong ML}]
\begin{itemize}
  \item \textbf{Mean, median, mode}: hiểu phân phối dữ liệu và phát hiện outlier.
  \item \textbf{Độ lệch chuẩn, phương sai}: đo độ biến thiên của dataset và phát hiện outlier.
  \item \textbf{Phân vị (percentiles)}: tóm tắt phân phối và nhận diện outlier.
  \item \textbf{Phân phối dữ liệu}: cách các điểm dữ liệu trải trên tập dữ liệu.
  \item \textbf{Độ lệch (skewness) và độ nhọn (kurtosis)}: đo hình dạng phân phối xác suất.
  \item \textbf{Bias và variance}: mô tả nguồn sai số trong dự đoán; overfitting/underfitting: Mục~\ref{sec:train-val-test}, bài \texttt{04}.
  \item \textbf{Kiểm định giả thuyết}: giả thuyết tạm thời có thể kiểm tra bằng dữ liệu.
  \item \textbf{Hồi quy tuyến tính}: thuật toán hồi quy supervised phổ biến nhất.
  \item \textbf{Hồi quy logistic}: thuật toán supervised quan trọng trong ML.
  \item \textbf{PCA}: chủ yếu dùng cho giảm chiều trong ML.
\end{itemize}
\end{phuongphap}
